\documentclass[11pt,a4paper]{article}

\usepackage[utf8]{inputenc}
\usepackage[T1]{fontenc}
\usepackage{amsmath,amssymb,bm}
\usepackage{booktabs}
\usepackage[margin=2.5cm]{geometry}
\usepackage{hyperref}
\usepackage{enumitem}

\hypersetup{
    colorlinks=true,
    linkcolor=blue,
    citecolor=blue,
    urlcolor=blue
}

\newcommand{\mb}[1]{\mathbf{#1}}
\newcommand{\zhat}{\hat{z}}
\newcommand{\Heff}{\mb{H}_{\text{eff}}}
\newcommand{\Hex}{\mb{H}_{\text{ex}}}
\newcommand{\Hdmi}{\mb{H}_{\text{DMI}}}
\newcommand{\Hanis}{\mb{H}_{\text{anis}}}
\newcommand{\Hdemag}{\mb{H}_{\text{demag}}}
\newcommand{\Hrkky}{\mb{H}_{\text{RKKY}}}
\newcommand{\Hext}{\mb{H}_{\text{ext}}}
\newcommand{\tCo}{t_{\text{Co}}}
\newcommand{\dRu}{d_{\text{Ru}}}
\newcommand{\Keff}{K_{\text{eff}}}
\newcommand{\Aex}{A_{\text{ex}}}
\newcommand{\Ms}{M_s}
\newcommand{\HK}{H_K}
\newcommand{\Dc}{D_c}
\newcommand{\Qpma}{Q_{\text{PMA}}}
\newcommand{\HRKKY}{H_{\text{RKKY}}}

\title{Phase Diagram of the SAF Skyrmion Simulator:\\
Theory, Reduced Units, and Implementation}
\author{Rui Barreira}
\date{May 2026}

\begin{document}
\maketitle

\begin{abstract}
This note documents the theoretical content of the phase-diagram
subsystem added to the SAF skyrmion simulator. The original
simulator (see \texttt{docs/theory.tex}) folded the thin-film
magnetostatic correction into the effective anisotropy
$\Keff = K - \mu_0 \Ms^2 / 2$. The phase-diagram subsystem
replaces that contraction with an explicit two-dimensional
Fourier-space demag, introduces a corrected total-energy
functional that respects the bilinear vs.\ external prefactor
convention, and classifies relaxed textures into eight phases
using order parameters that combine layer-averaged
magnetization, topological charge, an interlayer dot product,
and the angular content of the Fourier spectrum at the
dominant wavevector. We also adopt reduced units
$(D / \Dc,\; H_z / \HK)$ on the phase-diagram axes so the
boundaries are material-independent, and we extend the
initial-condition ensemble with a \emph{parallel} ferromagnet
seed so the sweep can resolve the spin-flop transition that
the antiparallel manifold alone cannot expose.
\end{abstract}

%======================================================================
\section{Why a New Energy Functional}
%======================================================================

\subsection{The $\Keff$ shortcut}

The base simulator computes the effective field with a single
\emph{local} anisotropy term whose prefactor is
$\Keff = K - \mu_0 \Ms^2 / 2$, absorbing the leading thin-film
magnetostatic contribution into the bulk anisotropy. The
identity that justifies this is, for a thin film of in-plane
extent $\gg \tCo$ in the uniform limit,
%
\begin{equation}\label{eq:keff_uniform}
\mathcal{H}_{\text{anis}} + \mathcal{H}_{\text{demag}}
    = -K \tCo \int m_z^2 \, dS
    + \frac{\mu_0 \Ms^2}{2} \, \tCo \int m_z^2 \, dS
    = -\Keff \, \tCo \int m_z^2 \, dS.
\end{equation}
%
For the uniform out-of-plane saturated state this is exact
and recovers the textbook shape anisotropy. For a textured
state with characteristic in-plane wavenumber $|\mb{k}| > 0$
the contraction \emph{fails}: the magnetostatic tensor in
Fourier space is not constant across $\mb{k}$, so the demag
contribution to a periodic texture differs from the uniform
prediction. Two consequences matter here:
%
\begin{enumerate}[nosep]
\item Textures with finite $|\mb{k}|$ generically gain energy
    relative to uniform FM, because the thin-film shape
    function is bounded by one:
    \begin{equation}\label{eq:shape_function}
    f(|\mb{k}| \tCo) = \frac{1 - e^{-|\mb{k}|\tCo}}
        {|\mb{k}|\tCo} \le 1,
    \end{equation}
    with $f \to 1$ at $|\mb{k}|\tCo \to 0$ (the uniform
    limit) and $f \to 0$ at $|\mb{k}|\tCo \to \infty$. The
    $\Keff$ contraction overestimates the demag penalty of a
    short-wavelength state by a factor $1/f$.
\item In a synthetic antiferromagnet the two layers couple
    magnetostatically across the Ru spacer. The interlayer
    contribution vanishes at $\mb{k} = 0$ (an infinite
    uniformly magnetized slab produces no external field) but
    is finite at $|\mb{k}| > 0$ with decay $e^{-|\mb{k}|\dRu}$.
    A local-anisotropy model has no access to this term.
\end{enumerate}

\subsection{Explicit 2D Fourier demag}

Each Co layer is treated as a slab of thickness $\tCo$
embedded in vacuum with periodic in-plane boundary conditions.
The magnetization on layer $\ell$ is decomposed in Fourier
modes $\widetilde{\mb{m}}^{(\ell)}(\mb{k})$, and the demag
field on layer $\ell$ from both layers reads
%
\begin{equation}\label{eq:demag_k}
\widetilde{\Hdemag}^{(\ell)}(\mb{k})
    = -\mu_0 \Ms \sum_{\ell'} \mb{N}^{(\ell, \ell')}(\mb{k})
        \cdot \widetilde{\mb{m}}^{(\ell')}(\mb{k}),
\end{equation}
%
with $\mb{N}^{(\ell, \ell')}(\mb{k})$ the dimensionless
$3 \times 3$ thin-film magnetostatic tensor. The diagonal
components used in the simulator are
%
\begin{align}
N_{xx}^{\text{self}}(\mb{k}) &= \frac{k_x^2}{k^2}
    \bigl(1 - f(k\tCo)\bigr), &
N_{yy}^{\text{self}}(\mb{k}) &= \frac{k_y^2}{k^2}
    \bigl(1 - f(k\tCo)\bigr), \nonumber\\
N_{zz}^{\text{self}}(\mb{k}) &= f(k\tCo), &
N_{xy}^{\text{self}}(\mb{k}) &= \frac{k_x k_y}{k^2}
    \bigl(1 - f(k\tCo)\bigr),
\label{eq:N_self}
\end{align}
%
with $k = |\mb{k}|$ and $f$ from~\eqref{eq:shape_function}.
The inter-layer tensor uses the same in-plane structure
multiplied by an exponentially decaying gap factor and an
overall sign flip on the out-of-plane component:
%
\begin{equation}\label{eq:N_inter}
\mb{N}^{(\text{top, bot})}(\mb{k})
    = S(k) \begin{pmatrix}
        k_x^2 / k^2 & k_x k_y / k^2 & 0 \\
        k_x k_y / k^2 & k_y^2 / k^2 & 0 \\
        0 & 0 & -1
    \end{pmatrix},
\qquad
S(k) = \frac{(1 - e^{-k\tCo})^2}{(k\tCo)^2}
    e^{-k\dRu},
\end{equation}
%
and vanishes at $\mb{k} = 0$, consistent with an infinite
parallel-plate slab producing no field outside itself. In
the limit $k\tCo \to 0$ on the self-tensor only,
\eqref{eq:N_self} reduces to $N_{zz}^{\text{self}}(0) = 1$ and
all in-plane components vanish, recovering the uniform
shape-anisotropy result
$(\Hdemag)_z = -\mu_0 \Ms m_z$ and
hence the $\Keff$ relation~\eqref{eq:keff_uniform}.

\subsection{Bare anisotropy at finite $\mb{k}$}

Because the demag now appears explicitly, the bulk
anisotropy prefactor used by the LLGS effective field is
the \emph{bare} $K$, not $\Keff$:
%
\begin{equation}
\Hanis^{(\ell)} = \frac{2 K_\ell}{\Ms} m_z^{(\ell)} \, \zhat,
\qquad K_\ell \in \{K_{\text{top}}, K_{\text{bot}}\}.
\end{equation}
%
The original simulator stores $\Keff$ in its anisotropy
prefactor, so the phase-diagram code recomputes
$2 K_\ell / \Ms$ from the raw material constants $(K_\ell,
\Ms)$ on each parameter namespace. The legacy field path
that uses $\Keff$ is left untouched so existing
current-driven analyses remain unaffected.

\subsection{Corrected total energy}

For a system with effective field expressed in tesla
($\mb{H}_T = \mu_0 \, \mb{H}_{A/m}$), the per-cell
self-energy
of a bilinear contribution $\mb{H}$ and the Zeeman term are
%
\begin{equation}\label{eq:E_density}
e_{\text{self}} = -\tfrac{1}{2}\, \Ms \,
    \mb{m} \cdot \mb{H}_T,
\qquad
e_Z = -\Ms \, \mb{m} \cdot \mb{H}_{\text{ext},T},
\end{equation}
%
without any additional $\mu_0$ prefactor because that
factor is already absorbed in the conversion from A/m to T.
The total energy summed over both Co layers therefore reads
%
\begin{equation}\label{eq:E_total}
\boxed{\;
E = -\tfrac{1}{2}\, V \Ms \!
    \sum_{\ell} \sum_{i}
    \mb{m}_i^{(\ell)} \cdot \mb{H}_{\text{int},i}^{(\ell)}
    - V \Ms \!
    \sum_{\ell} \sum_{i}
    \mb{m}_i^{(\ell)} \cdot \mb{H}_{\text{ext},T}
\;}
\end{equation}
%
with $V = \tCo a^2$ the per-site volume of one layer and
$\mb{H}_{\text{int}} =
\Hex + \Hdmi + \Hanis^{\text{bare}} + \Hrkky + \Hdemag$. The
$\tfrac{1}{2}$ prefactor prevents double counting of the
bilinear self-interactions; it is absent on the Zeeman term
because the external field is independent of the
magnetization. Substituting
$\Hanis^{\text{bare}} = (2K/\Ms) m_z \zhat$ recovers the
textbook density $-K m_z^2$, confirming the convention.

%======================================================================
\section{Reduced Units}
%======================================================================

A phase diagram in absolute $(D,\, H_z)$ collapses to a
single universal picture once the axes are normalized by the
material's natural scales. We adopt:
%
\begin{equation}\label{eq:reduced}
\boxed{\;
x \equiv \frac{D}{\Dc},
\qquad
y \equiv \frac{H_z}{\HK},
\;}
\end{equation}
%
with $\Dc$ the Bogdanov--Hubert critical DMI and $\HK$ the
anisotropy field. Both depend only on the material; they are
computed once per sweep and stored as scalars in the output.

\subsection{Critical DMI $\Dc$ (Bogdanov--Hubert)}

The uniform PMA ferromagnet is locally stable against
spontaneous helical winding only below
%
\begin{equation}\label{eq:Dc}
\Dc = \frac{4 \sqrt{\Aex \, \Keff}}{\pi},
\qquad
\Keff = \tfrac{1}{2}(K_{\text{top}} + K_{\text{bot}})
    - \tfrac{1}{2} \mu_0 \Ms^2.
\end{equation}
%
Above $\Dc$ the FM gains energy by canting into a
periodic helix; below it the FM is the ground state at
$H_z = 0$. The factor $4/\pi$ is the saddle-point result
for a sinusoidal spiral on an infinite plane with no
anisotropy in the easy-plane direction. We use the
layer-averaged $\Keff$ to fold the small asymmetry between
$K_{\text{top}}$ and $K_{\text{bot}}$ into a single scalar.
The formula is undefined for $\Keff \le 0$; the
implementation raises an exception in that easy-plane
regime, since no PMA FM exists to destabilize.

\subsection{Anisotropy field $\HK$}

The anisotropy field is the field at which the Zeeman
energy of a uniform out-of-plane ferromagnet equals the
anisotropy gap. In A/m units the textbook expression is
$\HK = 2 \Keff / (\mu_0 \Ms)$. The simulator's convention
keeps every field in tesla, so we store
%
\begin{equation}\label{eq:HK}
\HK^T = \mu_0 \HK = \frac{2 \Keff}{\Ms}
    \quad \text{(tesla)}.
\end{equation}

\subsection{PMA quality factor}

The dimensionless ratio
%
\begin{equation}\label{eq:Qpma}
\Qpma = \frac{\Keff}{\tfrac{1}{2}\mu_0 \Ms^2}
\end{equation}
%
sets which regime the material is in.
Table~\ref{tab:Qpma_regimes} summarizes the qualitative
behaviour for the SAF system; we obtain
$\Qpma \approx 0.007$ at $K = 1.294 \times 10^6$ J/m$^3$
(default) and $\Qpma \approx 0.245$ at
$K = 1.60 \times 10^6$ J/m$^3$ (the value used in this
work to access the chiral regime).
%
\begin{table}[h]
\centering
\small
\begin{tabular}{l p{5.0cm} p{5.0cm}}
\toprule
$\Qpma$ range & Ground state at $H_z = 0$ &
    Representative material \\
\midrule
$\lesssim 0.05$ & demag-driven bubbles / labyrinth
    & weakly-PMA Co (too thick) \\
$0.2$--$0.5$ & FM at large $|H_z|$;
    SkX/iSk in a window
    & Pt/Co/Ru/Co with engineered interface PMA \\
$> 1$ & uniform out-of-plane FM dominates
    & CoFeB / Heusler \\
\bottomrule
\end{tabular}
\caption{Qualitative PMA regimes set by the dimensionless
quality factor $\Qpma$ defined in~\eqref{eq:Qpma}. The
default simulator parameters fall in the first row; the
$K = 1.60$ MJ/m$^3$ choice used in the present sweep sits in
the middle row.}
\label{tab:Qpma_regimes}
\end{table}

\subsection{Material constants for the present sweep}

For $K = 1.60 \times 10^6$ J/m$^3$ in both layers and the
default $\Ms = 1.43$ MA/m, $\Aex = 16$ pJ/m, we obtain:
%
\begin{equation*}
\Keff \approx 0.315\ \text{MJ/m}^3,
\qquad
\Dc \approx 2.86\ \text{mJ/m}^2,
\qquad
\HK \approx 0.441\ \text{T},
\qquad
\Qpma \approx 0.245.
\end{equation*}
%
With $D$ swept over $[0, 6]$ mJ/m$^2$ and $H_z$ over
$[-0.5, +0.5]$ T, the reduced-unit coverage is
$x \in [0, 2.10]$ and $y \in [-1.13, +1.13]$, straddling
both the chiral threshold $x = 1$ and the SAF spin-flop
$|y| = \HRKKY / \HK \approx 0.47$.

%======================================================================
\section{Phase Classification}
%======================================================================

A relaxed configuration is labelled by a decision tree on
five order parameters. The labels are
$\{\text{FM\_anti}, \text{FM\_par+}, \text{FM\_par-},
\text{iSk}, \text{SkX}, \text{BX}, \text{SS},
\text{Lab}\}$ plus a sentinel \texttt{undetermined}.

\subsection{Order parameters}

For the top layer (the bottom layer is RKKY-coupled and
follows it up to canting):
%
\begin{itemize}[nosep]
\item Layer-averaged out-of-plane magnetization
    $\langle m_z \rangle$.
\item Layer dot product
    $\langle \mb{m}_{\text{top}} \cdot
    \mb{m}_{\text{bot}} \rangle$, the interlayer
    Néel order parameter. Cleanly antiparallel SAF
    states satisfy this $\approx -1$, parallel
    (spin-flopped) FM states satisfy $\approx +1$.
\item Topological charge of the top layer,
    $Q = \frac{1}{4\pi} \int \mb{m} \cdot
    (\partial_x \mb{m} \times \partial_y \mb{m}) \, dA$.
\item Power spectrum of $m_z(\mb{r})$. Two summary
    statistics are computed: a peak-to-background ratio
    on the radial spectrum (used to gate the
    ``periodic'' branch) and the
    2-fold and 6-fold angular harmonics of the azimuthal
    profile on a Gaussian ring of width
    $\sigma = 1.5\, \Delta k_{\text{grid}}$ around the
    dominant in-plane wavevector $k^\star$.
\item Topological charge per principal period,
    \begin{equation}\label{eq:q_per_period}
    q_p \equiv \frac{|Q|}{N_p},
    \qquad
    N_p \equiv
    \frac{k^\star L_x}{2\pi} \cdot
    \frac{k^\star L_y}{2\pi}.
    \end{equation}
    A hexagonal skyrmion lattice has $q_p \approx 1$
    (one full unit of charge per unit cell); a bubble
    lattice with the same FFT signature but trivial
    topology has $q_p \approx 0$.
\end{itemize}

\subsection{Decision tree}

Let $|\langle m_z \rangle| > 0.95$ mean ``uniform'',
$|\langle \mb{m}_{\text{top}} \cdot \mb{m}_{\text{bot}}
\rangle| > 0.9$ mean ``clean polar'', and
$P_6 / P_{\text{iso}} > 2$ with
$P_6 > P_2$ mean ``6-fold periodic'' (analogously for SS).
The classifier proceeds:
%
\begin{enumerate}[nosep]
\item If uniform: pick FM by interlayer dot product.
    Antiparallel ($\langle \mb{m}_{\text{top}} \cdot
    \mb{m}_{\text{bot}} \rangle < -0.9$) yields
    \textbf{FM\_anti}; parallel yields
    \textbf{FM\_par+} or \textbf{FM\_par-} based on the
    sign of $\langle m_z \rangle$. The uniform check is
    \emph{and}-coupled with the dot-product check so a
    textured antiparallel state with $\langle
    \mb{m}_{\text{top}} \cdot \mb{m}_{\text{bot}}
    \rangle \approx -1$ but
    $|\langle m_z \rangle| < 1$ falls through.
\item If $|Q| \in [0.5, 1.5]$ (one or a few isolated
    skyrmions, broad FFT): \textbf{iSk}.
\item If a 6-fold periodic ring is present:
    \textbf{SkX} when $q_p \ge 0.5$,
    \textbf{BX} otherwise.
\item If a 2-fold periodic ring is present:
    \textbf{SS} (spin spiral / helix).
\item If a periodic ring with no clean angular order:
    \textbf{Lab} (labyrinthine).
\item Fallbacks: $|Q| > 1.5$ without periodic ordering
    $\to$ \textbf{iSk} (multi-skyrmion family);
    $|Q| < 0.5$ without ordering $\to$ \textbf{Lab};
    else \textbf{undetermined}.
\end{enumerate}

\subsection{Why SkX vs.\ BX matters}

Both phases have 6-fold periodic $m_z(\mb{r})$ in the
FFT. They differ in whether each periodic minimum is
topologically protected. A real skyrmion lattice has
$|Q| = N_{\text{sk}}$ skyrmions per field of view, so
$q_p \to 1$. A bubble lattice (or a stripe / labyrinth
pattern with hexagonal symmetry) carries no winding and
$q_p \to 0$. At zero DMI no real SkX can exist by
construction; any 6-fold periodic ground state at
$D = 0$ is therefore BX. Conflating the two would
falsely report a chiral phase across regions of the
$(D, H_z)$ plane that are in fact demag-driven and
topologically trivial.

%======================================================================
\section{Initial-Condition Ensemble}
%======================================================================

At each $(D, H_z)$ grid point the relaxation is run from
seven initial conditions; the lowest-energy converged
result is taken as the ground state. The ensemble is
designed to seed every basin we know to be physically
relevant:
%
\begin{itemize}[nosep]
\item Three independent random spheres
    ($\varphi \sim U[0, 2\pi]$, $\cos\theta \sim U[-1, 1]$,
    different seeds) to explore the texture manifold
    without geometric bias. The bottom layer is set to
    $-\mb{m}_{\text{top}}$ so the RKKY antiparallel
    constraint is respected initially.
\item \texttt{fm\_anti}: antiparallel SAF
    ($\mb{m}_{\text{top}} = \pm\zhat$,
    $\mb{m}_{\text{bot}} = \mp\zhat$). Tracks the FM
    branch of the antiparallel manifold.
\item \texttt{fm\_par}: parallel FM
    ($\mb{m}_{\text{top}} = \mb{m}_{\text{bot}} =
    \text{sign}(H_z) \zhat$). Tracks the FM branch of
    the parallel manifold, accessible above the
    spin-flop. Required to expose the
    $H_z$-driven antiparallel $\to$ parallel transition;
    without it the $(D, H_z)$ diagram is symmetric under
    $H_z \to -H_z$ up to numerical noise.
\item \texttt{saf\_skyrmion}: a Néel skyrmion in the top
    layer with antiparallel pair below.
\item \texttt{stripe}: a helical stripe of pitch
    $\lambda = 4\pi \Aex / D$ (the natural DMI helix
    period). Floored at $4 a$ to remain resolvable; for
    $D = 0$ the period is set to the largest available
    helical seed.
\end{itemize}

\subsection{Zero-Zeeman SAF and the spin-flop}

Inside the antiparallel manifold ($\mb{m}_{\text{top}} =
-\mb{m}_{\text{bot}}$ cell by cell) the Zeeman density is
%
\begin{equation}
e_Z = -\Ms (\mb{m}_{\text{top}} +
    \mb{m}_{\text{bot}}) \cdot \mb{H}_{\text{ext}, T}
    \equiv 0,
\end{equation}
%
identically, regardless of $|H_z|$. The applied field is
decoupled from every state inside this manifold, including
FM\_anti, BX, SkX, iSk, SS, Lab. A genuine $H_z$-driven
phase boundary therefore requires the parallel manifold,
which becomes energetically favorable above the spin-flop:
%
\begin{equation}\label{eq:spinflop}
|H_z| \gtrsim \HRKKY \approx 0.205\ \text{T},
\end{equation}
%
where the Zeeman gain from aligning both layers exceeds
the RKKY cost of breaking antiparallel coupling. The
\texttt{fm\_par} seed allows the relaxation to find this
basin; without it the sweep would only sample
antiparallel-manifold states and the high-$|y|$ regions
of the phase diagram would be incorrectly attributed to
the antiparallel FM rather than to the parallel FM.

%======================================================================
\section{Numerical Workflow}
%======================================================================

\subsection{Sweep}

The sweep iterates over a Cartesian grid of $(D, H_z)$
values; for each grid point and each of the seven ICs the
relaxation runs an RK4 LLGS integrator with $J = 0$ and
optional over-damped Gilbert factor $\alpha_{\text{relax}}
= 1$ until both convergence criteria fire:
%
\begin{equation}\label{eq:conv}
\max_i \bigl| \mb{m}_i \times
    (\mb{m}_i \times \Heff(\mb{r}_i)) \bigr|
    < \tau_{\text{tol}},
\qquad
\left| \frac{\Delta E}{E} \right| < \epsilon_E,
\end{equation}
%
checked every $10^3$ steps (defaults
$\tau_{\text{tol}} = 10^{-5}$ T, $\epsilon_E = 10^{-8}$,
safety cutoff $2 \times 10^5$ steps). The
tangential-torque criterion measures
$|\mb{m} \times (\mb{m} \times \Heff)|$ rather than
$|\dot{\mb{m}}|$ to avoid contamination from any
$|\mb{m}| \ne 1$ numerical drift.

The lowest-energy converged IC at each grid point is taken
as the ground state, and per-IC observables plus the
ground-state texture are written to a single compressed
NPZ.

\subsection{Reduced-unit plotting}

The plotter defaults to $(D / \Dc, H_z / \HK)$ axes. The
critical-DMI vertical line $x = 1$ and the spin-flop
$|y| = \HRKKY / \HK$ are drawn on the phase map for
reference. Absolute units (mJ/m$^2$, T) are available as a
fallback.

\subsection{Classifier verification}

The classifier is verified on synthetic textures whose
labels are known by construction:
%
\begin{itemize}[nosep]
\item Uniform $(\pm\zhat, \mp\zhat)$ $\to$ FM\_anti.
\item Uniform $(\pm\zhat, \pm\zhat)$ $\to$ FM\_par$\pm$.
\item Helical stripe at the natural pitch $\to$ SS.
\item Sum of three plane waves at $0$, $60$, $120$
    degrees with Néel winding $\to$ SkX (with $q_p$
    above threshold) or BX (without).
\end{itemize}

%======================================================================
\section{Summary}
%======================================================================

The phase-diagram subsystem replaces the $\Keff$ thin-film
contraction with an explicit 2D Fourier-space demag
\eqref{eq:demag_k}--\eqref{eq:N_inter}, restores the bare
anisotropy prefactor at the LLGS effective field, and writes
total energies through the bilinear-vs.-external convention
\eqref{eq:E_density}--\eqref{eq:E_total}. Phase diagrams are
expressed in the dimensionless coordinates
$(D/\Dc, H_z/\HK)$ defined by~\eqref{eq:Dc}--\eqref{eq:HK},
classified by a decision tree that distinguishes
topologically non-trivial skyrmion lattices from trivial
bubble lattices via the per-period charge
$q_p$~\eqref{eq:q_per_period}, and built from a 7-member IC
ensemble that includes a parallel-FM seed so the spin-flop
transition above $|H_z| \gtrsim \HRKKY$ is resolved. The
choice $K = 1.60$ MJ/m$^3$ places the material at
$\Qpma \approx 0.25$, in the regime where chiral phases are
expected; the absolute sweep ranges $D \in [0, 6]$
mJ/m$^2$ and $H_z \in [-0.5, +0.5]$ T translate to
$D/\Dc \in [0, 2.10]$ and $H_z/\HK \in [-1.13, +1.13]$.

%======================================================================
\begin{thebibliography}{1}
\bibitem{pham2024}
V.~T.~Pham, N.~Sisodia, I.~Di~Manici,
J.~Urrestarazu-Larra\~{n}aga et al.,
``Fast current-induced skyrmion motion in synthetic
antiferromagnets,''
\emph{Nature Physics} (2024).
\bibitem{bogdanov1994}
A.~N.~Bogdanov and A.~Hubert,
``Thermodynamically stable magnetic vortex states in
magnetic crystals,''
\emph{Journal of Magnetism and Magnetic Materials}
\textbf{138}, 255--269 (1994).
\bibitem{sampaio2013}
J.~Sampaio, V.~Cros, S.~Rohart, A.~Thiaville, A.~Fert,
``Nucleation, stability and current-induced motion of
isolated magnetic skyrmions in nanostructures,''
\emph{Nature Nanotechnology} \textbf{8}, 839 (2013).
\bibitem{rohart2013}
S.~Rohart and A.~Thiaville,
``Skyrmion confinement in ultrathin film nanostructures
in the presence of Dzyaloshinskii-Moriya interaction,''
\emph{Physical Review B} \textbf{88}, 184422 (2013).
\bibitem{buttner2018}
F.~B\"{u}ttner, I.~Lemesh, and G.~S.~D.~Beach,
``Theory of isolated magnetic skyrmions: From
fundamentals to room temperature applications,''
\emph{Scientific Reports} \textbf{8}, 4464 (2018).
\bibitem{bernandmantel2020}
A.~Bernand-Mantel, C.~B.~Muratov, and T.~M.~Simon,
``Unraveling the role of dipolar versus
Dzyaloshinskii-Moriya interactions in stabilizing
compact magnetic skyrmions,''
\emph{Physical Review B} \textbf{101}, 045416 (2020).
\end{thebibliography}

\end{document}
